% texte.tex — f174v, vv. 1001–1043 (conversion initiale des travaux antérieurs)
\vers[n=1002]{ΜΝ.}{Mnésiloque}{ὦ τοξόθ᾽ ἱκετεύω σε.}{Archer, je t'en supplie...}
\vers[suite]{ΤΟ.}{Archer}{μή μ᾽ ἰκετεῦσι σύ.}{Ne me supplie pas.}
\vers[n=1003]{ΜΝ.}{Mnésiloque}{χάλασον τὸν ἧλον.}{Desserre la cheville.}
\vers[suite]{ΤΟ.}{Archer}{ἀλλὰ ταῦτα δρᾶσ᾽ ἐγώ.}{C'est ce que je vais faire}
\vers[n=1004]{ΜΝ.}{Mnésiloque}{οἴμοι κακοδαίμων, μᾶλλον ἐπικρούεις σύ γε.}{Aïe, misère! tu la pousses davantage.}
\vers[n=1005]{ΤΟ.}{Archer}{ἔτι μᾶλλο βοῦλις;}{Encore plus, veux-tu?}
\vers[suite]{ΜΝ.}{Mnésiloque}{ἀτταταῖ ἰατταταῖ·}{Aie, aïe, aïe, aie, aïe, aie !}
\vers[n=1006]{}{}{κακῶς ἀπόλοιο.}{la peste soit de toi.}
\vers[suite]{ΤΟ.}{Archer}{σῖγα κακοδαίμων γέρον.}{Silence, misérable vieillard.}
\vers[n=1007]{}{}{πέρ᾽ ἐγὼ 'ξινίγκι πορμός, ἴνα πυλάξι σοι.}{Allons, je vais, moi, apporter une natte, pour faire ma garde. \didasc{Il sort.}}
\vers[n=1008]{ΜΝ.}{Mnésiloque}{ταυτὶ τὰ βέλτιστ᾽ ἀπολέλαυκ᾽ Εὐριπίδου.}{Ces douceurs-là j'en ai l'obligation à Euripide.}
\vers[n=1009]{}{}{ἔα· θεοί, Ζεῦ σῶτερ, εἰσὶν ἐλπίδες.}{\didasc{Il aperçoit dans la coulisse Euripide costumé en Persée.} Hein ! dieux ! Zeus sauveur ! il y a espoir.}
\vers[n=1010]{}{}{ἁνὴρ ἔοικεν οὐ προδώσειν, ἀλλά μοι}{L'homme ne semble pas vouloir m'abandonner. Mais il m'a}
\vers[n=1011]{}{}{σημεῖον ὑπεδήλωσε Περσεὺς ἐκδραμών,}{donné à entendre, en apparaissant en Persée,}
\vers[n=1012]{}{}{ὅτι δεῖ με γίγνεσθ᾽ Ἀνδρομέδαν· πάντως δέ μοι}{qu'il me faut devenir Andromède. Aussi bien,}
\vers[n=1013]{}{}{τὰ δέσμ᾽ ὑπάρχει. δῆλον οὖν τοῦτ᾽ ἔσθ᾽ ὅτι}{les chaînes, j'en suis muni. Il est donc clair}
\vers[n=1015]{}{}{ἥξει με σώσων· οὐ γὰρ ἂν παρέπτετο.}{qu'il viendra me délivrer. Autrement il n'aurait pas dirigé son vol de ce côté.}
\vers[n=1016, indent=2em, findent=1.2em, metre={lyr.}]{}{}{φίλαι παρθένοι φίλαι,}{Vierges, vierges amies, puissé-je}
\vers[n=1017, indent=2em, findent=1.2em]{}{}{πῶς ἂν ἀπέλθοιμι καὶ}{partir d'ici et}
\vers[n=1018, indent=2em, findent=1.2em]{}{}{τὸν Σκύθην λάθοιμι;}{tromper le Scythe !}
\vers[n=1019, indent=2em, findent=1.2em]{}{}{κλύεις; ὦ προσᾴδουσ' ἀϋτὰς ἐν ἄντροις,}{M'entends-tu, ô toi qui en des antres répètes mes chants ?}
\vers[n=1020, indent=2em, findent=1.2em]{}{}{κατάνευσον, ἔασον ὡς}{acquiesce, souffre que}
\vers[n=1021, indent=2em, findent=1.2em]{}{}{τὴν γυναῖκά μ᾽ ἐλθεῖν.}{j'aille retrouver ma femme.}
\vers[n=1022, indent=2em, findent=1.2em]{}{}{ἄνοικτος ὅς μ᾽ ἔδησε τὸν}{Sans merci est celui qui m'enchaîna,}
\vers[n=1023, indent=2em, findent=1.2em]{}{}{πολυπονώτατον βροτῶν·}{des mortels le plus éprouvé.}
\vers[n=1024, indent=2em, findent=1.2em]{}{}{μόλις δὲ γραῖαν ἀποφυγὼν}{Échappé à grand'peine à une vieille}
\vers[n=1025, indent=2em, findent=1.2em]{}{}{σαπρὰν ἀπωλόμην ὅμως.}{guenipe, je ne laissai pas d'être perdu.}
\vers[n=1026, indent=2em, findent=1.2em]{}{}{ὅδε γὰρ ὁ Σκύθης φύλαξ πάλαι ἐφέστηκ᾽· ὀλοὸν}{Car depuis un temps ce Scythe se tient en faction : il m'a suspendu, brisé, sans amis,}
\vers[n=1027, indent=2em, findent=1.2em]{}{}{ἄφιλον ἐκρέμασε κόραξι δεῖπνον.}{en pâture aux corbeaux. Vois-tu, ce n'est pas}
\vers[n=1030, indent=2em, findent=1.2em]{}{}{ὁρᾷς; οὐ χοροῖσιν οὐδ᾽ ὑφ᾽ ἡλίκων νεανίδων}{pour des chœurs de danse}
\vers[n=1031, indent=2em, findent=1.2em]{}{}{ψῆφον ἕστηκ᾽ ἔχουσα κημόν,}{ni parmi des jeunes filles de mon âge}
\vers[n=1032, indent=2em, findent=1.2em]{}{}{ἀλλ᾽ ἐν πυκνοῖς δεσμοῖσιν ἐμπεπλεγμένη}{que je suis ici debout avec cornet et jeton.}
\vers[n=1033, indent=2em, findent=1.2em]{}{}{κήτει βορὰ Γλαυκέτῃ πρόκειμαι,}{Mais enchaînée en d'étroits liens,}
\vers[n=1034, indent=2em, findent=1.2em]{}{}{γαμηλίῳ μὲν οὐ ξὺν}{je me vois exposée à la voracité du monstre Glaucétès.}
\vers[n=1035, indent=2em, findent=1.2em]{}{}{παιῶνι δεσμίῳ δὲ}{Lamentez sur moi un péan, non d'hymen, mais de chaînes,}
\vers[n=1036, indent=2em, findent=1.2em]{}{}{γοᾶσθέ μ᾽ ὦ γυναῖκες ὡς}{ô femmes, car,}
\vers[n=1037, indent=2em, findent=1.2em]{}{}{μέλεα μὲν πέπονθα μέλεος,}{pitoyable, j'ai pitoyablement souffert.}
\vers[n=1038, indent=2em, findent=1.2em]{}{}{ὢ τάλας ἐγὼ τάλας,}{Malheureux que je suis, malheureux !}
\vers[n=1039, indent=2em, findent=1.2em]{}{}{ἀπὸ δὲ συγγόνων ἀλλ᾽ ἄνομα πάθεα,}{Et entre d'autres tortures impies infligées par mes proches,}
\vers[n=1040, indent=2em, findent=1.2em]{}{}{†φῶτά λιτομένα†,}{implorant l'homme}
\vers[n=1040, indent=2em, findent=1.2em]{}{}{πολυδάκρυτον Ἀίδα γόον ψεφαυγοῦς.}{d'une lamentation pleine de larmes de l'obscur Hadès.}
\vers[n=1042, indent=2em, findent=1.2em]{}{}{αἰαῖ αἰαῖ ἒ ἒ,}{ah ah! ah ah!}
\vers[n=1043, indent=2em, findent=1.2em]{}{}{ὃς ἔμ᾽ ἀπεξύρησε πρῶτον,}{l'homme qui me rasa d'abord,}

\vers[n=1044, indent=2em, findent=1.2em]{}{}{ὃς ἐμὲ κροκόεντ᾽ ἐνέδυσεν·}{me revêtit de cette crocote,}
